\frfilename{mt1.tex}
\versiondate{30.9.02}
\copyrightdate{1994}

\def\volumename{Minimum yang tak tereduksi}

\newvolume{1}


Dalam jilid pengantar ini saya menguraikan, pada tingkat yang saya harap
sesuai bagi mahasiswa yang belum memiliki pengetahuan sebelumnya tentang
integral Lebesgue (atau bahkan integral Riemann) dan yang hanya memiliki
bekal dasar (namun tuntas) dalam teknik analisis
$\epsilon$-$\delta$, teori ukuran dan integrasi hingga
teorema-teorema konvergensi (\S123). Saya menambahkan bab ketiga (Bab
13) yang memuat berbagai hasil tambahan, sebagian besar dipilih karena
merupakan materi relatif elementer yang diperlukan untuk topik-topik yang
dibahas dalam Jilid 2 tetapi tidak memiliki tempat yang wajar di sana.

Judul jilid ini sedikit lebih tegas daripada yang ingin saya coba
pertanggungjawabkan {\it au pied de la lettre}. Saya tentu akan
menyebut konstruksi ukuran Lebesgue pada $\Bbb R$ (\S114),
definisi integral pada ruang ukur abstrak (\S122), dan
teorema-teorema konvergensi (\S123) sebagai hal-hal yang mutlak perlu.
Namun, seorang pengajar yang ingin segera beralih ke topik-topik lebih
lanjut akan mendapati bahwa sebagian besar Bab 13 dapat dikesampingkan
untuk sementara. Di sini saya mengatakan `pengajar', bukan
`mahasiswa', karena jika Anda belajar sendiri, saya pikir Anda
sebaiknya berusaha maju lebih lambat daripada yang dituntut teks, bukan
lebih cepat; menurut pandangan saya, gagasan-gagasan ini benar-benar
sulit, dan saya pikir Anda sebaiknya meluangkan waktu untuk berlatih
sebanyak mungkin pada tingkat yang relatif elementer.

Barangkali inilah saat yang tepat untuk mengemukakan beberapa pemikiran
umum tentang pengajaran teori ukuran. Kini saya telah mengajar analisis
selama lebih dari tiga puluh tahun, dan salah satu dari sedikit hal yang
tetap sama sepanjang masa itu adalah perasaan, yang hampir universal di
kalangan pengajar analisis, bahwa kita tidak melayani sebagian besar
mahasiswa kita dengan baik. Kita semua pernah menjumpai mahasiswa yang
tidak bodoh -- bahkan sesungguhnya cukup pandai dalam matematika --
tetapi tampaknya mengalami kesulitan yang tidak sebanding dalam analisis
yang ketat. Mereka begitu letih dan kehilangan semangat akibat persoalan
teknis sehingga tidak mampu memahami atau bahkan menggunakan pengetahuan
yang telah mereka peroleh. Reaksi yang wajar terhadap keadaan ini adalah
mencoba membuat mata kuliah lebih singkat dan lebih mudah. Namun, saya
pikir hal itu justru makin memperbesar kemungkinan bahwa pada akhir
semester mahasiswa Anda akan terdampar di semak berduri di tengah
perjalanan mendaki gunung. Khususnya dalam hal ukuran Lebesgue, Anda
terancam menghabiskan dua puluh jam untuk mengajari mereka cara
mengintegralkan fungsi karakteristik himpunan bilangan rasional. Bukan untuk
itulah pokok bahasan ini ada. Penyajian Lebesgue sendiri mengenai pokok
bahasan ini ({\smc Lebesgue 1904}, {\smc Lebesgue 1918}) menekankan
teorema-teorema konvergensi dan Teorema Dasar Kalkulus. Yang pertama telah
saya tempatkan dalam Jilid 1 dan yang terakhir dalam Jilid 2, tetapi menurut
saya, kecuali mahasiswa Anda sendiri ingin tahu kapan kita dapat berharap
bahwa limit dan integral bisa dipertukarkan, atau
fungsi mana yang merupakan integral tak tentu, atau seperti apa
ruang-ruang hasil pelengkapan $C([0,1])$ di bawah norma
$\|\,\|_1$,  $\|\,\|_2$, maka
akan sangat sulit bagi mereka untuk memperoleh sesuatu dari materi ini;
dan jika Anda benar-benar tidak dapat mencapai tahap menjelaskan
sekurang-kurangnya dua dari perkara ini dengan istilah yang dapat mereka
pahami, mungkin tidak ada gunanya memulai. Saya sendiri lebih memilih
menghilangkan cukup banyak pembuktian daripada baru sampai pada
teorema-teorema yang menjadi alasan diciptakannya pokok bahasan ini
sedemikian terlambat dalam mata kuliah sehingga dua pertiga kelas saya
sudah menyerah sebelum teorema-teorema itu dibahas.

Tentu saja, saya dan orang-orang lain juga pernah menempuh jalan itu,
dengan hasil yang tidak lebih baik (meskipun biasanya dengan mahasiswa
yang lebih bahagia) daripada hasil yang kita peroleh dengan menggarap setiap
{\it perincian} dan menuntaskan setiap {\it langkah} dalam pembuktian. Hampir
setiap kali saya dimintai pendapat oleh seorang
nonspesialis yang ingin diberi tahu sebuah teorema yang akan memecahkan
masalahnya, saya kembali diingatkan bahwa matematika murni pada umumnya,
dan analisis pada khususnya, tidak terletak dalam {\it teorema-teorema},
melainkan dalam {\it pembuktian-pembuktian}. Sejauh saya berhasil
menjawab pertanyaan-pertanyaan semacam itu, biasanya hal itu terjadi
dengan melakukan penyesuaian kecil pada sebuah argumen baku untuk
menghasilkan sebuah teorema tak baku.
Gagasan-gagasannya terletak dalam perinciannya. Anda belum memahami
konstruksi \Caratheodory{} (\S113) sampai Anda setidak-tidaknya dapat
mereproduksi secara andal argumen yang menunjukkan bahwa konstruksi itu
berhasil. Pada akhirnya, tidak ada pilihan selain menelusuri setiap
langkah, dan meskipun sesekali saya tanpa ampun membuang topik-topik yang
menurut saya hanya bersifat marginal, saya telah berusaha memastikan --
yang sering kali harus dibayar dengan kepedantisan -- bahwa setiap
gagasan yang diperlukan telah ditandai.

Oleh karena itu, ketika menghadapi kelas tertentu, saya percaya bahwa
seorang pengajar harus berkompromi antara keluasan cakupan dan
kelengkapan. Persis kompromi mana yang paling sesuai akan bergantung pada
faktor-faktor yang hanya akan membuang waktu jika saya mencoba
menebaknya. Jilid ini dimaksudkan sebagai salah satu teks yang mungkin
digunakan sebagai landasan mata kuliah; tetapi saya berharap tidak ada
dosen yang menugasi kelasnya membacanya sekian halaman setiap minggu.
Tujuan utama saya adalah menyediakan landasan yang ringkas dan koheren
untuk mendirikan struktur jilid-jilid selanjutnya. Untuk itu, pada beberapa
bagian saya harus mengikuti pendekatan yang menempuh jalur
sedikit lebih sulit demi mengembangkan teknik yang lebih halus. (Mungkin
yang paling menonjol di antaranya adalah kegigihan saya dalam menegaskan bahwa
suatu fungsi terintegralkan tidak harus terdefinisi di setiap titik pada
ruang ukur yang mendasarinya; lihat \S\S121-122.)
Setiap pengajar bertanggung jawab memutuskan sendiri apakah
penyempurnaan-penyempurnaan semacam itu sesuai dengan kebutuhan
mahasiswanya, dan jika tidak, menunjukkan kepada mereka penyesuaian apa
yang diperlukan.

Paragraf-paragraf di atas ditujukan kepada para pengajar yang, menurut
anggapan, menguasai pokok bahasan ini -- tentu saja melampaui tingkat
yang dibahas dalam jilid ini -- dan memiliki akses ke sejumlah buku
yang sangat baik dan sudah tersedia, sehingga jika mereka bersusah payah
memikirkan tujuan mereka, mereka semestinya dapat memilih unsur-unsur
penyajian saya yang sesuai. Namun, saya juga harus mempertimbangkan
keadaan seorang mahasiswa yang mulai mempelajari materi ini sendiri.
Saya percaya bahwa dari apa yang telah saya tulis, Anda sudah memahami
bahwa Anda tidak perlu takut melihat ke depan. Bahkan, ada baiknya Anda
beralih ke Jilid 2 dan mengkaji salah satu teorema menakjubkan di
sana -- misalnya Teorema Dasar Kalkulus (\S222), atau, jika Anda sangat
berambisi, hukum kuat bilangan besar (\S273) -- lalu menggunakan indeks
dan rujukan silang untuk mencoba memperoleh pembuktiannya dari
prinsip-prinsip pertama. Jika berhasil, Anda sepenuhnya berhak memberi
selamat kepada diri sendiri. Namun, pada masa-masa ketika keberhasilan
terasa sukar diraih, Anda sebaiknya mengerjakan secara sistematis
`latihan-latihan dasar' dalam bagian-bagian yang tampaknya relevan; dan
jika segala upaya lain gagal, mulailah lagi dari awal. Matematika adalah
pokok bahasan yang sulit, itulah sebabnya matematika layak ditekuni, dan
hampir setiap bagian di sini memuat suatu gagasan esensial yang tidak
mungkin Anda temukan sendiri.

\frnewpage
